\chapter{Aprendizaje de máquina: predecir}
\label{cap:s-predecir}
\textit{Materia: Aprendizaje de máquina. Pregunta: ¿cómo estará cada lugar el mes que entra, y cuánta gente llegará en
los próximos 12 meses?}

\section{Qué significa que una computadora ``aprenda''}
Un modelo de aprendizaje de máquina es una receta que se ajusta sola con ejemplos del pasado. Se le muestran muchos casos
(``este mes el índice era tal y el mes siguiente el lugar estuvo concurrido'') y busca la regla que mejor los explica. Después
se le da un caso nuevo y responde con esa regla.

El proyecto tiene \textbf{dos} problemas de predicción distintos, y conviene no confundirlos:
\begin{center}
\small
\begin{tabularx}{\linewidth}{lYY}
\encabezado & \h{El Radar} & \h{El Pronóstico} \\
\textbf{Qué predice} & Una \emph{palabra}: tranquilo, concurrido o saturado & Un \emph{número}: cuántos visitantes \\
\textbf{Para cuándo} & El mes siguiente & De 1 a 12 meses \\
\textbf{Contra qué se compara} & ``Igual que este mes'' & ``Lo mismo que el mismo mes del año pasado'' \\
\bottomrule
\end{tabularx}
\normalsize
\end{center}

\begin{porque}
\textbf{Siempre contra una regla simple.} Un modelo complicado que no le gana a una regla de sentido común no sirve. Por eso
cada modelo se compara contra la regla más simple posible, y se dice cuánto le gana (o si no le gana).
\end{porque}

\section{El Radar: ¿cómo estará el lugar el mes que entra?}
\textbf{Qué ve el modelo}: el índice de presión de este mes y de los dos anteriores, cuánto cambió, la época del año y si es
uno de los cinco lugares. \textbf{Qué responde}: la probabilidad de cada estado el mes siguiente.

Se compararon tres modelos:
\begin{itemize}
\item \textbf{Regresión logística}: le da a cada estado un puntaje que suma lo que empuja cada dato, y convierte los
  puntajes en probabilidades que suman 100\,\%.
\item \textbf{Random Forest}: 400 árboles de decisión, cada uno entrenado con una muestra distinta; el resultado es el
  promedio de sus votos.
\item \textbf{Gradient Boosting}: árboles pequeños en fila, donde cada uno corrige los errores del anterior.
\end{itemize}

\begin{ejemplo}[ · Mahahual, agosto de 2026]
Mahahual tuvo en julio un índice de 0.731, muy alto por los cruceros. La regresión logística calcula un puntaje para cada
estado y los convierte en probabilidades:
\begin{itemize}
\item concurrido: \textbf{60.1\,\%};
\item saturado: \textbf{39.9\,\%};
\item tranquilo: casi 0.
\end{itemize}
La página dice: Mahahual seguirá concurrido, con 4 de 10 probabilidades de saturarse. (Mahahual no se promueve; sirve para
ver que el modelo detecta presión donde la hay.)
\end{ejemplo}

\subsection{Cómo se probó: viajar al pasado}
\begin{intuicion}
La forma escolar de probar un modelo es apartar al azar una parte de los datos. Aquí sería trampa: el modelo podría
aprender con agosto y probarse con julio, es decir, \emph{ver el futuro}. Lo honesto es ``viajar al pasado'': pararse en
agosto de 2025, entrenar solo con lo que se sabía hasta entonces, predecir septiembre y comparar con lo que pasó. Luego
pararse en septiembre, y así hasta julio de 2026.
\end{intuicion}
Fueron 12 meses de prueba, 156 predicciones y 30 cambios reales de estado.

\begin{center}
\small
\begin{tabular}{lrrr}
\encabezado \h{Modelo} & \h{Aciertos (de 156)} & \h{Cambios anticipados (de 30)} & \h{Falsas alarmas} \\
``Igual que este mes'' & 126 & 0 & 0 \\
\textbf{Regresión logística} & \textbf{130} & \textbf{8} & 4 \\
Random Forest & 130 & 6 & 2 \\
Gradient Boosting & 128 & 7 & 5 \\
\bottomrule
\end{tabular}
\normalsize
\end{center}

\textbf{Criterio de Brandon}: el que más acierta y, si empatan, el que más cambios anticipa, porque la regla anti-colapso
necesita \emph{anticipar}. Ganó la regresión logística.

\begin{teprofe}
\emph{``Tu modelo apenas le gana a la regla simple, ¿para qué sirve?''} Es cierto, y se dice: el estado se repite 8 de cada
10 meses, así que ``igual que este mes'' acierta mucho. El modelo gana 4 aciertos, pero su valor está en otra parte:
anticipa 8 de los 30 cambios, y la regla simple no anticipa ninguno. Además, en los cinco lugares casi no hubo cambios
(2 en 48 casos), así que \textbf{no se puede afirmar que el modelo anticipe saturaciones en el sur}; anticipa cambios en los
lugares de referencia. Decirlo así es más fuerte que exagerar.
\end{teprofe}

\section{El Pronóstico: ¿cuánta gente llegará?}
\subsection{Qué se pronostica y qué no}
Solo se pronostica lo que está \textbf{medido} y llega hasta 2026:
\begin{itemize}
\item visitantes a Oxtankah (Bahía) y a las tres zonas de la Ruta, desde 2016 (INAH);
\item cruces desde Belice por Chetumal, desde 2019 (gobierno del estado);
\item la ocupación de Cancún y de la Riviera Maya, como referencia.
\end{itemize}
Maya Ka'an y la Laguna Milagros no tienen ninguna serie de visitantes: no se pronostican, y se dice.

\begin{porque}
\textbf{Los meses cerrados no enseñan.} Si una zona estuvo cerrada por obras y marcó 0 visitantes, un modelo que aprenda de
ese cero creería que en ese mes puede no llegar nadie ``por temporada''. Por eso los meses de cierre, los meses a medias
(justo antes de cerrar o después de reabrir) y la pandemia se marcan y \textbf{no entrenan}, aunque su valor se conserva.
Además, las zonas reabrieron con menos gente que antes (Kohunlich tuvo 21,850 visitantes en 2025 contra 42,813 en 2019),
así que el nivel se toma de los meses después de la reapertura.
\end{porque}

\subsection{Cinco formas de predecir}
\begin{enumerate}
\item \textbf{La regla simple}: lo mismo que el mismo mes del año pasado.
\item \textbf{Holt-Winters sin tendencia}: se le quita la temporada a la serie, se sigue su ``nivel de fondo'' mezclando
  cada dato nuevo con lo que ya se creía, y se le vuelve a poner la temporada.
\item \textbf{Holt-Winters con tendencia}: igual, más una pendiente que se va apagando.
\item \textbf{Regresión con clima}: el nivel de cada etapa (antes de la pandemia, después, desde la reapertura), el efecto de
  cada mes, la lluvia fuera de lo normal y las tormentas.
\item \textbf{Gradient Boosting}: árboles que aprenden de los últimos meses y del mismo mes del año anterior.
\end{enumerate}

\begin{ejemplo}[ · la Ruta en enero de 2027, con Holt-Winters]
El ``nivel de fondo'' de la Ruta (sin temporada) era de unos 5,332 visitantes al mes. En julio de 2026 llegaron 5,057, y
como julio es un mes casi normal (0.953), sin temporada eso equivale a 5,307. El modelo cree un 11.8\,\% en el dato nuevo y
un 88.2\,\% en lo que ya sabía:
\[
0.118\times5{,}307+0.882\times5{,}332\approx5{,}329 .
\]
Para enero se le vuelve a poner la temporada: $5{,}329\times1.606=\mathbf{8{,}559}$ visitantes.
\end{ejemplo}

\begin{ejemplo}[ · la Bahía en diciembre de 2026, con la regresión]
Desde que Oxtankah reabrió en 2024, un enero típico tiene unos 1,139 visitantes. Diciembre trae 15.1\,\% más que enero, y
nunca ha habido tormenta en diciembre. Pronóstico: $1{,}139\times1.151\approx\mathbf{1{,}311}$ visitantes.
\end{ejemplo}

\subsection{El rango: ``muy probablemente entre tanto y tanto''}
Un pronóstico solo, sin rango, engaña: parece más seguro de lo que es. Cada pronóstico lleva un \textbf{rango del 90\,\%}.

\begin{intuicion}
Se revisan los errores que el modelo cometió en el pasado en pronósticos parecidos. Si 9 de cada 10 fueron menores que
cierto tamaño, lo razonable es decir que el próximo tampoco se equivocará más que eso. No se supone ninguna curva: se usan
los errores reales. Este método se llama \emph{conformal}.
\end{intuicion}

\begin{ejemplo}[ · otra vez la Bahía en diciembre de 2026]
De 105 errores pasados de pronósticos a 4–6 meses, el que deja abajo al 90\,\% equivale a equivocarse unos 37\,\% hacia arriba
o 27\,\% hacia abajo. Entonces:
\begin{itemize}
\item mínimo: $1{,}311\times0.731=959$;
\item máximo: $1{,}311\times1.368=1{,}793$.
\end{itemize}
Se esperan \textbf{1,311 visitantes, muy probablemente entre 959 y 1,793}. Diciembre de 2025 tuvo 1,224: dentro del rango.
\end{ejemplo}

\textbf{¿Se cumple de verdad?} Se midió en el pasado cuántas veces el dato real cayó dentro del rango, y se reporta aunque
salga mal: en la Bahía, 91 de cada 100 veces; en Belice, 94; en la Ruta, solo \textbf{80}.

\subsection{Cuál se eligió}
\textbf{Regla de Brandon}: ``menor error con rango $\geq$ 80\,\%''. De los modelos cuyo rango se cumple al menos 80 de cada
100 veces, el que menos se equivoca.

\begin{center}
\small
\begin{tabular}{llrr}
\encabezado \h{Lugar} & \h{Modelo elegido} & \h{Error típico} & \h{Rango cumplido} \\
Bahía & Regresión con clima & 15.7\,\% & 91.3\,\% \\
Ruta & Regresión con clima & 21.6\,\% & 80.3\,\% \\
Chetumal (Belice) & Regresión con clima & 11.0\,\% & 94.1\,\% \\
Cancún (referencia) & La regla simple & 3.8\,\% & 89.3\,\% \\
\bottomrule
\end{tabular}
\normalsize
\end{center}

\begin{porque}
\textbf{Por qué no el de menor error, sin más.} En la Ruta, el de menor error (Holt-Winters sin tendencia) tiene un rango
que solo atrapa el dato real 72 de cada 100 veces: un rango ``del 90\,\%'' que falla tanto engaña al que planea.
\textbf{Por qué no un solo modelo para todo.} En Cancún, la regresión se equivoca 58\,\% más que repetir el año anterior;
la ocupación de Cancún cambia muy poco de un año a otro.
\end{porque}

\subsection{Lo que se declara}
\begin{itemize}
\item \textbf{La Ruta falló en 2023}: las visitas cayeron 10.7\,\% sin ningún aviso en la historia, y ese año el rango solo
  se cumplió 67 de cada 100 veces. En 2025–2026 vuelve a cumplirse.
\item \textbf{El clima casi no mejora el pronóstico}: sin lluvia ni tormentas, la regresión se equivoca casi igual. Lo que
  funciona es separar las etapas y el efecto de cada mes. Aun así, el clima sí se asocia: en la Bahía, cada 100 mm de
  lluvia sobre lo normal van con 10.9\,\% menos visitantes.
\item \textbf{La tendencia engaña en tramos cortos}: Holt-Winters con tendencia llegó a pronosticar $-240$ visitantes
  (¡negativo!) en la Bahía, porque aprendió la ``tendencia'' en plena reapertura.
\end{itemize}

\textbf{Para la campaña}: el Pronóstico define la temporada alta de cada lugar y cuánta gente se espera; con eso el
modelo del dinero decide en qué meses anunciar (capítulo~\ref{cap:s-dinero}).
